%!TEX program = pdflatex
% ─────────────────────────────────────────────────────────────
%  Lebenslauf — Friedrich Wilke Grosche
%  Standalone pdfLaTeX, kompilierbar auf Overleaf
% ─────────────────────────────────────────────────────────────
\documentclass[a4paper,10pt]{article}

% ── Packages ──────────────────────────────────────────────────
\usepackage[T1]{fontenc}
\usepackage[utf8]{inputenc}
\usepackage[ngerman]{babel}
\usepackage[
  left=0pt, right=0pt, top=0pt, bottom=0pt,
  nohead, nofoot, nomarginpar
]{geometry}
\usepackage{tikz}
\usetikzlibrary{calc,positioning}
\usepackage{xcolor}
\usepackage{fontawesome5}
\usepackage{enumitem}
\usepackage{hyperref}
\usepackage{graphicx}
\usepackage{parskip}
\usepackage{tabularx}
\usepackage{setspace}
\usepackage{amssymb}

\hypersetup{
  colorlinks=true,
  urlcolor=sidebarAccent,
  linkcolor=sidebarAccent,
  pdfauthor={Friedrich Wilke Grosche},
  pdftitle={Lebenslauf -- Friedrich Wilke Grosche},
}

\pagestyle{empty}

% ── Colour palette ────────────────────────────────────────────
\definecolor{sidebar}{HTML}{4B5C6B}        % dark blue-grey sidebar
\definecolor{sidebarText}{HTML}{FFFFFF}     % white text on sidebar
\definecolor{sidebarAccent}{HTML}{3E98A0}   % teal accent for headings
\definecolor{accentGold}{HTML}{C8A951}      % gold underlines
\definecolor{mainHeading}{HTML}{1B2A41}     % dark navy for main headings
\definecolor{mainText}{HTML}{333333}        % body text
\definecolor{mainLight}{HTML}{777777}       % lighter text (dates, org)
\definecolor{ruleColor}{HTML}{1B2A41}       % horizontal rules

% ── Layout constants ──────────────────────────────────────────
\newlength{\sidebarwidth}
\setlength{\sidebarwidth}{6.8cm}
\newlength{\mainwidth}
\setlength{\mainwidth}{\dimexpr\paperwidth-\sidebarwidth\relax}
\newlength{\sidebarpadding}
\setlength{\sidebarpadding}{0.6cm}
\newlength{\mainpadding}
\setlength{\mainpadding}{0.8cm}

% ── Reusable commands ─────────────────────────────────────────
% Sidebar section heading (teal + gold underline)
\newcommand{\sidebarsection}[1]{%
  \vspace{10pt}%
  {\color{sidebarAccent}\bfseries\large\MakeUppercase{#1}}\\[-2pt]%
  {\color{accentGold}\rule{3cm}{1.2pt}}%
  \vspace{6pt}\par%
}

% Main section heading (navy with lines above/below)
\newcommand{\mainsection}[1]{%
  \vspace{14pt}%
  {\color{ruleColor}\hrule height 0.8pt}%
  \vspace{6pt}%
  {\color{mainHeading}\bfseries\large\MakeUppercase{#1}}%
  \vspace{4pt}%
  {\color{ruleColor}\hrule height 0.8pt}%
  \vspace{8pt}\par%
}

% Decorative diamond divider
\newcommand{\divider}{%
  \begin{center}%
    \vspace{2pt}%
    {\color{mainLight}\small$\diamond$\hspace{6pt}$\diamond$\hspace{6pt}$\diamond$}%
    \vspace{2pt}%
  \end{center}%
}

% Experience entry
\newcommand{\expentry}[4]{%
  % #1 = title, #2 = dates, #3 = org, #4 = description
  {\color{mainHeading}\bfseries #1}\hfill{\color{mainLight}\small\itshape #2}\\[1pt]%
  {\color{sidebarAccent}\small\itshape #3}\\[4pt]%
  {\color{mainText}\small #4}\par%
  \vspace{6pt}%
}

% Education entry (sidebar)
\newcommand{\eduentry}[4]{%
  % #1 = title, #2 = dates, #3 = institution, #4 = city
  {\color{sidebarText}\bfseries\small\MakeUppercase{#1}}\\[1pt]%
  {\color{sidebarText}\small #2}\\[1pt]%
  {\color{sidebarText}\small #3}\\[0pt]%
  {\color{sidebarText}\small #4}\par%
  \vspace{8pt}%
}

% Sidebar check-item
\newcommand{\checkitem}[1]{%
  {\color{sidebarAccent}\faCheck}~~{\color{sidebarText}\small #1}\par\vspace{3pt}%
}

% Sidebar bullet item
\newcommand{\sideitem}[1]{%
  {\color{sidebarAccent}\scriptsize\raisebox{1pt}{$\blacksquare$}}~~{\color{sidebarText}\small #1}\par\vspace{2pt}%
}


% ══════════════════════════════════════════════════════════════
\begin{document}

% ══════════════  PAGE 1  ══════════════════════════════════════
\noindent
\begin{tikzpicture}[remember picture, overlay]
  % Sidebar background
  \fill[sidebar] (current page.north west) rectangle
    ([xshift=\sidebarwidth]current page.south west);
\end{tikzpicture}%
%
\noindent
\begin{minipage}[t][\textheight]{\sidebarwidth}
  \begin{center}
    \vspace{0.7cm}

    % ── Photo placeholder ──
    % Ersetze 'wilkeedit.png' durch deinen Dateinamen in Overleaf
    % \includegraphics[width=4cm,height=4cm]{wilkeedit.png}
    \begin{tikzpicture}
      \draw[sidebarAccent, line width=2pt] (0,0) circle (2cm);
      \node[text=sidebarText, font=\small] at (0,0) {Foto hier einf\"ugen};
    \end{tikzpicture}

    \vspace{0.5cm}
  \end{center}

  \hspace{\sidebarpadding}%
  \begin{minipage}[t]{\dimexpr\sidebarwidth-2\sidebarpadding\relax}

    % ── KONTAKT ──
    \sidebarsection{Kontakt}
    {\color{sidebarText}\small
      \faIcon{map-marker-alt}~~Z\"urich, Schweiz\\[3pt]
      \faIcon{phone-alt}~~+41 76 265 81 51\\[3pt]
      \faIcon{envelope}~~\href{mailto:wilkegrosche@gmail.com}{wilkegrosche@gmail.com}\\[3pt]
      \faIcon{github}~~\href{https://github.com/wgrosche}{github.com/wgrosche}\\[3pt]
      \faIcon{linkedin}~~\href{https://www.linkedin.com/in/wilke/}{linkedin.com/in/wilke}%
    }

    % ── EXPERTISE ──
    \sidebarsection{Expertise}
    \checkitem{Python (15+ Jahre, Fortgeschritten)}
    \checkitem{C++ (Kompetent)}
    \checkitem{CasADi \textbullet{} IPOPT}
    \checkitem{PyTorch \textbullet{} TensorFlow}
    \checkitem{OpenCV \textbullet{} MMDetection}
    \checkitem{Docker \textbullet{} Slurm \textbullet{} RunAI}
    \checkitem{Git \textbullet{} \LaTeX{} \textbullet{} CMake}
    \checkitem{Django \textbullet{} NumPy \textbullet{} SciPy}

    % ── SPRACHEN ──
    \sidebarsection{Sprachen}
    \sideitem{Englisch (Muttersprache)}
    \sideitem{Deutsch (Muttersprache)}
    \sideitem{Franz\"osisch (fliessend)}

    % ── SOFT SKILLS ──
    \sidebarsection{Soft Skills}
    \sideitem{Problemabstraktion \& Formalisierung}
    \sideitem{Interdisziplin\"are Kommunikation}
    \sideitem{Technische Pr\"asentation}

  \end{minipage}
\end{minipage}%
%
% ── MAIN COLUMN (PAGE 1) ─────────────────────────────────────
\begin{minipage}[t][\textheight]{\mainwidth}
  \hspace{\mainpadding}%
  \begin{minipage}[t]{\dimexpr\mainwidth-2\mainpadding\relax}
    \vspace{0.7cm}

    % ── NAME / TITLE ──
    \begin{center}
      {\color{mainHeading}\fontsize{28}{32}\selectfont\bfseries
       F\,R\,I\,E\,D\,R\,I\,C\,H~~\,W\,I\,L\,K\,E~~\,G\,R\,O\,S\,C\,H\,E}\\[6pt]
      {\color{mainLight}\normalsize Research Engineer}
    \end{center}
    \divider

    % ── ÜBER MICH ──
    \mainsection{\"Uber mich}
    {\color{mainText}\small
      Deutsch-britischer Physik-Absolvent (MSc, ETH Z\"urich) mit 3 Jahren Erfahrung als Research
      Engineer am EPFL CVLab. Schwerpunkt auf der Verbindung von Regelungstheorie, Computer
      Vision und maschinellem Lernen. Experte in der Entwicklung hochperformanter
      Autonomie-Stacks f\"ur schnelle Starrfl\"ugelflugzeuge (100\,m/s) mit Kernkompetenz in der
      \"Uberf\"uhrung komplexer physikalischer Systeme in handhabbare Optimierungsprobleme.
    }

    % ── BERUFSERFAHRUNG ──
    \mainsection{Berufserfahrung}

    {\color{mainHeading}\bfseries Research Engineer}\hfill{\color{mainLight}\small\itshape 09.2022 -- 09.2025}\\[1pt]
    {\color{sidebarAccent}\small\itshape EPFL Computer Vision Lab (CVLab), Lausanne, Schweiz}\\[6pt]

    {\color{mainText}\small
      \textbf{Zeitoptimale MPC f\"ur Starrfl\"ugel-Drohnen}
      \textit{(Kooperation mit dem Schweizerischen Forschungszentrum)} ---
      Entwicklung zeitoptimaler MPC und Trajektorienoptimierung f\"ur eine Starrfl\"ugel-Drohne
      unter F5B-Wettbewerbsbedingungen (100\,m/s, ca.\ 50 Runden \`a 300\,m in 3\,min). Aufbau
      eines differenzierbaren aerodynamischen Surrogatmodells unter Integration von
      CFD-Daten und physikalischen Grundmodellen.
      (\href{https://github.com/wgrosche/AIrcraft}{github.com/wgrosche/AIrcraft}).
      (2{,}5 Jahre)
    }
    \vspace{6pt}

    {\color{mainText}\small
      \textbf{Kameralokalisierung und -kalibrierung f\"ur Multi-View-Tracking} ---
      Leitung der Entwicklung einer Multi-Kamera-Lokalisierungs- und Kalibrierungssuite;
      automatisierte Achsenausrichtung mittels RANSAC-Ebenenanpassung und Kamerah\"ohen-Priors
      zur zuverl\"assigen Weltkoordinaten-Rekonstruktion.
      (\href{https://github.com/cvlab-epfl/MARMOT}{github.com/cvlab-epfl/MARMOT}). (1 Jahr)
    }
    \vspace{6pt}

    {\color{mainText}\small
      \textbf{ML-Modelltraining f\"ur Multi-View-Personentracking} ---
      Datenerweiterung, Pipelining und Training eines Graph Neural Networks f\"ur
      Langzeit-Personentracking mit Verdeckungen (modifizierte ByteTrack- und
      YoloX-Versionen). (1 Jahr)
    }
    \vspace{6pt}

    {\color{mainText}\small
      \textbf{Annotationspipeline f\"ur Multi-View-Datensatz} ---
      Leitung der Entwicklung eines webbasierten Multi-View-Annotationstools
      (\href{https://github.com/cvlab-epfl/multicam-gt}{github.com/cvlab-epfl/multicam-gt})
      f\"ur den \"offentlichen SCOUT-Datensatz
      (\href{https://scout.epfl.ch/}{scout.epfl.ch});
      automatisiertes Tracklet-Merging reduzierte die manuelle Annotationszeit um ca.\ 90\,\%.
      (8 Monate)
    }
    \vspace{6pt}

    {\color{mainText}\small
      \textbf{Physikbasierte Herz- und Kathetersimulation} ---
      Entwicklung einer physikbasierten Katheter-Herz-Interaktionssimulation f\"ur virtuelle
      Herzkatheterisierung (Cosserat-Stabmodell, SDF-basierte Kollisionserkennung) in Echtzeit.
      (5 Monate)
    }

  \end{minipage}
\end{minipage}

% ══════════════  PAGE 2  ══════════════════════════════════════
\newpage
\noindent
\begin{tikzpicture}[remember picture, overlay]
  \fill[sidebar] (current page.north west) rectangle
    ([xshift=\sidebarwidth]current page.south west);
\end{tikzpicture}%
%
\noindent
\begin{minipage}[t][\textheight]{\sidebarwidth}
  \hspace{\sidebarpadding}%
  \begin{minipage}[t]{\dimexpr\sidebarwidth-2\sidebarpadding\relax}
    \vspace{0.7cm}

    % ── AUSBILDUNG ──
    \sidebarsection{Ausbildung}

    \eduentry{M.Sc.\ Physik}{09.2019 -- 04.2022}{ETH Z\"urich}{Schweiz}
    {\color{sidebarText}\scriptsize\itshape Thesis: Model-Agnostic Explainability\\for ML Regression in Physical Systems\\
    Betreuer: Dr.\ A.\ Adelmann (PSI),\\ Prof.\ Dr.\ Ce Zhang (ETH Z\"urich)}
    \vspace{10pt}

    \eduentry{B.Sc.\ Physik}{09.2015 -- 09.2019}{ETH Z\"urich}{Schweiz}
    {\color{sidebarText}\scriptsize\itshape Thesis: Optimising Flux Gate Positions\\for Active Magnetic Shielding\\
    Betreuer: Prof.\ Dr.\ K.\,S.\ Kirch (ETH Z\"urich)}
    \vspace{10pt}

    \eduentry{A-Levels}{09.2013 -- 09.2015}{Hills Road Sixth Form College}{Cambridge, GB}
    {\color{sidebarText}\scriptsize\itshape Physics, Maths, Further Maths,\\Chemistry, German}
    \vspace{4pt}

    % ── PERSÖNLICH ──
    \sidebarsection{Pers\"onlich}
    \sideitem{geb.\ 1996}
    \sideitem{Deutsch-britisch}
    \sideitem{Ledig}

    % ── INTERESSEN ──
    \sidebarsection{Interessen}
    \sideitem{Kompetitive Brettspiele (Eurogames)}
    \sideitem{Fantasy-Literatur}
    \sideitem{Wandern, Bouldern, Skifahren}
    \sideitem{Kreuzwortr\"atsel}

    % ── REFERENZEN ──
    \sidebarsection{Referenzen}
    {\color{sidebarText}\small
      \textbf{Prof.\ Pascal Fua}\\
      {\scriptsize Leiter CVLab, EPFL}\\[6pt]
      \textbf{Dr.\ Andreas Adelmann}\\
      {\scriptsize Masterthesis-Betreuer, PSI}\\[6pt]
      \textbf{Dr.\ Martin Engilberge}\\
      {\scriptsize Projektleiter Tracking}\\[4pt]
      {\scriptsize\itshape (Kontaktdaten auf Anfrage)}
    }

  \end{minipage}
\end{minipage}%
%
% ── MAIN COLUMN (PAGE 2) ─────────────────────────────────────
\begin{minipage}[t][\textheight]{\mainwidth}
  \hspace{\mainpadding}%
  \begin{minipage}[t]{\dimexpr\mainwidth-2\mainpadding\relax}
    \vspace{0.7cm}

    % ── Repeated name header ──
    \begin{center}
      {\color{mainHeading}\fontsize{28}{32}\selectfont\bfseries
       F\,R\,I\,E\,D\,R\,I\,C\,H~~\,W\,I\,L\,K\,E~~\,G\,R\,O\,S\,C\,H\,E}\\[6pt]
      {\color{mainLight}\normalsize Research Engineer}
    \end{center}
    \divider

    % ── TECHNISCHES DETAIL ──
    \mainsection{Technisches Detail: Projekte}

    {\color{mainText}\small
      \textbf{Zeitoptimale MPC f\"ur Starrfl\"ugel-Drohnen}
      \textit{(Kooperation mit Microsoft Research)} ---
      Trajektorienoptimierung und zeitoptimale MPC f\"ur eine unmotorisierte
      Starrfl\"ugel-Drohne f\"ur den F5B-Wettbewerb (bis zu 100\,m/s, ca.\ 50 Runden
      \`a 300\,m in 3\,min), einschliesslich aerodynamischer Systemidentifikation aus
      Windkanal-, CFD- und Flugdaten. Aufbau eines differenzierbaren Surrogatmodells
      mittels CasADi und IPOPT, Reduktion der aerodynamischen Kraftberechnung von
      ca.\ 8\,Stunden (CFD) auf ca.\ 1\,ms. Eigenst\"andige Herleitung einer
      Drehratenkorrrektur zur Ber\"ucksichtigung dynamischer Effekte, die in der CFD
      fehlen --- L\"osung von Zeitintegrationsinstabilit\"aten im Trajektorienverlauf.
    }
    \vspace{8pt}

    {\color{mainText}\small
      \textbf{Kameralokalisierung und -kalibrierung f\"ur Multi-View-Tracking} ---
      Leitung der Kameralokalisierung und -kalibrierung f\"ur ein Multi-View-Trackingsystem,
      einschliesslich einer automatisierten Achsenausrichtungsmethode f\"ur
      OpenSFM-Rekonstruktionen mittels RANSAC-Ebenenanpassung auf Keypoint-dichten
      Regionen in Kombination mit Kamerah\"ohen-Priors zur Bestimmung der $z$-Achsenorientierung.
      Anpassung von ByteTrack und YOLOX f\"ur Bodenebenen-Betrieb durch Projektion der
      Detektionen auf den Boden zur sinnvollen kamera\"ubergreifenden Korrespondenz.
    }
    \vspace{8pt}

    {\color{mainText}\small
      \textbf{Annotationspipeline f\"ur Multi-View-Datensatz} ---
      Aufbau einer Annotationspipeline f\"ur einen Multi-View-Personentracking-Datensatz
      mit zusammenh\"angend \"uberlappenden Sichtfeldern \"uber 25 Kameras
      (\href{https://scout.epfl.ch/}{scout.epfl.ch}), die Trajektorienanalyse \"uber
      ca.\ 10-min\"utige Zeitr\"aume erm\"oglicht. Automatisiertes n\"ahebasiertes
      Tracklet-Merging und kamera\"ubergreifende Korrespondenz reduzierten die manuelle
      Annotationszeit um ca.\ 90\,\%.
    }

    % ── PUBLIKATIONEN ──
    \mainsection{Publikationen}
    {\color{mainText}\small
      \begin{enumerate}[leftmargin=1.4em, itemsep=4pt, label={[\arabic*]}]
        \item M.\ Engilberge, \textbf{F.\ Wilke Grosche} und P.\ Fua.
              \textit{No Identity, no problem: Motion through detection for people tracking.}
              2024. arXiv: \href{https://arxiv.org/abs/2411.16466}{2411.16466}.
        \item M.\ Engilberge et al.\
              \textit{Unified People Tracking with Graph Neural Networks.}
              2025. arXiv: \href{https://arxiv.org/abs/2507.08494}{2507.08494}.
      \end{enumerate}
    }

  \end{minipage}
\end{minipage}

\end{document}
